\documentclass[11pt,a4paper]{amsart}

\usepackage[margin=1.1in]{geometry}
\usepackage{amsmath,amssymb,amsthm,mathtools}
\usepackage{mathrsfs}
\usepackage{stmaryrd}
\usepackage{enumitem}
\usepackage{booktabs}
\usepackage{xcolor}
\usepackage[colorlinks=true,linkcolor=black!70!blue,citecolor=black!70!blue,urlcolor=black!70!blue]{hyperref}
\usepackage{microtype}

\theoremstyle{plain}
\newtheorem{theorem}{Theorem}[section]
\newtheorem{proposition}[theorem]{Proposition}
\newtheorem{lemma}[theorem]{Lemma}
\newtheorem{corollary}[theorem]{Corollary}

\theoremstyle{definition}
\newtheorem{definition}[theorem]{Definition}
\newtheorem{notation}[theorem]{Notation}
\newtheorem{example}[theorem]{Example}

\theoremstyle{remark}
\newtheorem{remark}[theorem]{Remark}

\DeclareMathOperator{\Adj}{Adj}
\DeclareMathOperator{\Ker}{ker}
\DeclareMathOperator{\Img}{im}
\DeclareMathOperator{\diag}{diag}
\DeclareMathOperator{\rank}{rank}
\DeclareMathOperator{\tr}{tr}
\DeclareMathOperator{\id}{id}
\DeclareMathOperator{\End}{End}
\DeclareMathOperator{\Aut}{Aut}
\DeclareMathOperator{\sgn}{sgn}
\newcommand{\RR}{\mathbb{R}}
\newcommand{\ZZ}{\mathbb{Z}}
\newcommand{\NN}{\mathbb{N}}
\newcommand{\FF}{\mathbb{F}}
\newcommand{\II}{\mathbb{I}}
\newcommand{\CC}{\mathbb{C}}
\newcommand{\HH}{\mathbb{H}}
\newcommand{\cC}{\mathcal{C}}
\newcommand{\cE}{\mathcal{E}}
\newcommand{\cF}{\mathcal{F}}
\newcommand{\cV}{\mathcal{V}}
\newcommand{\cL}{\mathcal{L}}
\newcommand{\cG}{\mathcal{G}}
\newcommand{\cN}{\mathcal{N}}
\newcommand{\kronprod}{\otimes_{\!K}}
\newcommand{\sheet}{\text{\rm sh}}
\newcommand{\wrap}{\operatorname{wrap}}
\newcommand{\tG}{\widetilde{G}}
\newcommand{\tV}{\widetilde{V}}
\newcommand{\tL}{\widetilde{L}}
\newcommand{\Lscalar}{L^{\operatorname{sc}}}
\newcommand{\Lsigned}{L^{\operatorname{sg}}}
\newcommand{\comp}{\pi_0}
\newcommand{\pair}[2]{\langle #1, #2 \rangle}
\newcommand{\defeq}{\mathrel{\mathop:}=}

\begin{document}

\title[Kernel Dimension of the $\ZZ/2$ Connection Laplacian]{%
Kernel Dimension of the $\ZZ/2$ Connection Laplacian \\
on a Finite Simple Graph \\
{\small (A Formalisation in Lean~4 / Mathlib v4.11.0)}}

\author{J.~Vilela}
\address{Independent}
\email{jesusvilela@outlook.com}
\date{\today}

\subjclass[2020]{05C50, 05C22, 55N31, 68V15}
\keywords{signed graph, connection Laplacian, double cover, balanced
  component, formal proof, Lean, Mathlib}

\begin{abstract}
We present a fully formalised, machine-checked proof library in
Lean~4 / Mathlib~v4.11.0 for the kernel-dimension theory of the
$\ZZ/2$ \emph{connection Laplacian} of a finite simple graph~$G$.
A connection graph is a finite simple graph equipped with a
designated set $W\subseteq E(G)$ of \emph{wrap edges}; on each
edge~$e$ we attach a $2\times 2$ real matrix (the Pauli~$\sigma_x$
on wrap edges, $\II_2$ otherwise), yielding a block Laplacian
$L_{\operatorname{Mob}}$.
Our principal results, each proved from the three standard Lean
axioms (\texttt{propext}, \texttt{Classical.choice},
\texttt{Quot.sound}) and nothing else, are the following.

\begin{enumerate}[label=\textbf{(\roman*)}]
  \item A Hadamard--Kronecker conjugation decomposes
    $L_{\operatorname{Mob}}$ into the block-diagonal sum of the
    ordinary (\emph{scalar}) Laplacian and a \emph{signed} Laplacian,
    both on~$G$.
  \item The orientation double cover $\tG$ of~$G$ inherits a trivial
    connection, and its scalar Laplacian splits similarly.
  \item A base connected component $C\in \comp(G)$ is
    \emph{balanced} iff the induced map $\comp(\tG)\to \comp(G)$ has
    fibre of cardinality~$2$ over~$C$.
  \item The kernel dimension of the signed Laplacian equals the number
    of balanced components; the full $\ZZ/2$-connection Laplacian has
    kernel dimension $\# \comp(G) + \#\{\text{balanced components}\}$
    (Mobius) and $2\,\#\comp(G)$ (flat).
  \item On the $n$-cycle the Fourier basis diagonalises both modes
    explicitly.
\end{enumerate}

The bespoke arguments are the Hadamard rotation of part~(i)--(ii)
(a Kronecker-conjugation computation) and the double-cover lift
of part~(iii) (a walk induction). Along the way we observe and
correct an incorrect statement that appeared in our pre-registered
plan: the identity
$\dim\Ker L^{\operatorname{sg}}(G) = \#\comp(G) -
  \#\{\text{components of odd total wrap}\}$
holds only on $\ZZ/2$-cycles, and must in general be replaced by the
cohomological formula $\dim\Ker L^{\operatorname{sg}}(G) =
\#\{\text{balanced components}\}$, which is the content of
Theorem~\ref{thm:signed-kernel}. The Lean sources and a numerical
cross-check on every simple graph of order~$\le 5$
($3{,}456$ configurations, zero disagreements)~--~and an extended
fuzzer sweep at orders $n\le 7$ ($456{,}950$ configurations, also zero
disagreements) accompany this paper.
\end{abstract}

\maketitle

\tableofcontents

%=============================================================
\section{Introduction}
%=============================================================

Let $G=(V,E)$ be a finite simple graph and let $W\subseteq E$ be
a distinguished set of \emph{wrap} edges. Attach to each
$e\in E$ a $2\times 2$ real matrix $\Phi(e)$ by
\[
  \Phi(e) \defeq \begin{cases} \sigma_x & \text{if } e\in W,\\
    \II_2 & \text{if } e\notin W, \end{cases}
\qquad
  \sigma_x = \begin{pmatrix}0&1\\1&0\end{pmatrix}.
\]
The \emph{$\ZZ/2$-connection Laplacian} of $(G,W)$ is the
$V\times V$ block matrix with $2\times 2$ blocks
\begin{equation}\label{eq:connection-laplacian}
  L_{\operatorname{Mob}}[u,v] \defeq
    \begin{cases}
      \deg(u)\,\II_2 & u=v,\\
      -\,\Phi(\{u,v\}) & \{u,v\}\in E,\\
      0 & \text{otherwise.}
    \end{cases}
\end{equation}
Equivalently, $L_{\operatorname{Mob}} = D\otimes \II_2 -
\sum_{e\in E} \Phi(e)\otimes \mathbf{1}_e$, where $\mathbf{1}_e$ is
the $V\times V$ edge-indicator of~$e$.
The \emph{flat} Laplacian $L_{\operatorname{flat}}$ is obtained by
the same recipe with $\Phi\equiv \II_2$, i.e. by ignoring the wrap
structure; it agrees with two copies of the ordinary graph Laplacian.
The guiding question answered in this paper is:

\begin{center}
  \emph{In terms of $(G,W)$, what is $\dim\Ker L_{\operatorname{Mob}}$
  and when is it strictly smaller than $\dim\Ker L_{\operatorname{flat}}$?}
\end{center}

This is the finite-dimensional, discrete analogue of the spectrum of
the Laplacian on a flat $\ZZ/2$ line bundle over a graph: the
kernel detects global sections of the bundle, and the flat bundle
always has two independent sections per component.
The Mobius bundle admits a non-vanishing section over a component
exactly when the \emph{wrap} holonomy of the bundle around every
cycle in that component is trivial; this is the condition of
\emph{balance} in the language of Harary, Cartwright and Zaslavsky
\cite{HararyCartwright, Zaslavsky82}.

\subsection*{Organisation and results}

Section~\ref{sec:prelim} sets the notation and fixes the
Lean-formalised data. Section~\ref{sec:decomp} proves the
Hadamard--Kronecker decomposition that reduces the block
Laplacian to two ordinary Laplacians (one scalar, one signed).
Section~\ref{sec:cover} constructs the orientation double cover
and proves that its scalar Laplacian splits as the direct sum of
the scalar and signed Laplacians of~$G$. Section~\ref{sec:coho}
formalises the balance notion as the vanishing of a
$\ZZ/2$-cohomology class and
computes the fibre cardinality of the induced map on
$\pi_0$. Section~\ref{sec:recognition} assembles everything into
the \emph{recognition theorems} which count kernel dimension by
balanced components. Section~\ref{sec:cycles} specialises to the
$n$-cycle, where the Fourier basis diagonalises both
Laplacians. Section~\ref{sec:formalisation} summarises the
formal Lean development.

Throughout, every displayed equation and every numbered theorem
is machine-checked in Lean~4 against Mathlib~v4.11.0. The
complete proof scripts and an exhaustive numerical verification
are included in the accompanying Lean project.

%=============================================================
\section{Preliminaries}\label{sec:prelim}
%=============================================================

\subsection{Connection graphs}

A \emph{connection graph} is a quadruple $G=(V,E,\Adj,W)$ where
$V$ is a finite type, $\Adj\subseteq V\times V$ defines a
(mathematically simple) graph on~$V$, $E=\{\{u,v\}:\Adj(u,v)\}$
is its edge set and $W\subseteq E$ is an arbitrary subset, the
\emph{wrap} edges. In Lean we keep the wrap data as a predicate
\texttt{wrap : G.graph.edgeSet $\to$ Prop} and carry explicit
\texttt{Fintype}/\texttt{DecidableEq}/\texttt{DecidableRel}
instances throughout.

\subsection{Pauli and Hadamard}

The $2\times 2$ Pauli matrix $\sigma_x=\begin{pmatrix}0&1\\1&0\end{pmatrix}$
is involutive, symmetric and has eigenvalues~$\pm 1$, with
symmetric eigenvector $(1,1)^\top$ and antisymmetric
$(1,-1)^\top$. The \emph{Hadamard rotation}
\begin{equation}\label{eq:Rhat}
  \widehat R \defeq \begin{pmatrix} 1 & 1 \\ 1 & -1 \end{pmatrix},
  \qquad \widehat R\,\widehat R = 2\,\II_2,
\end{equation}
simultaneously diagonalises $\sigma_x$ and $\II_2$ in the sense
that $\widehat R\,\sigma_x\,\widehat R = 2\,\diag(1,-1)$ and
$\widehat R\,\II_2\,\widehat R = 2\,\II_2$. In the Lean
formalisation we keep two copies of~$\widehat R$: one indexed by
$\{0,1\}=\operatorname{Fin}2$ (for the connection-block
decomposition of Section~\ref{sec:decomp}) and one indexed by
$\operatorname{Bool}$ (for the cover-sheet decomposition of
Section~\ref{sec:cover}). They are the same matrix under the
obvious bijection $\operatorname{Bool}\cong \operatorname{Fin}2$.

\subsection{Scalar and signed Laplacians}

Let $G$ be a connection graph. The ordinary \emph{scalar}
Laplacian $\Lscalar \in \RR^{V\times V}$ is
\[
  \Lscalar[u,v] \defeq
    \begin{cases} \deg(u) & u=v,\\
      -1 & \{u,v\}\in E,\\ 0 & \text{otherwise,}\end{cases}
\]
i.e.\ Mathlib's
\texttt{SimpleGraph.lapMatrix}.
The \emph{signed Mobius} Laplacian $\Lsigned \in \RR^{V\times V}$
encodes the wrap sign on the off-diagonal:
\[
  \Lsigned[u,v] \defeq
    \begin{cases} \deg(u) & u=v,\\
      +1 & \{u,v\}\in E,\ \{u,v\}\in W,\\
      -1 & \{u,v\}\in E,\ \{u,v\}\notin W,\\
      0 & \text{otherwise.}\end{cases}
\]
Both are real symmetric and positive semidefinite as quadratic forms
on $\RR^V$. For $\Lscalar$ this is Mathlib's
\texttt{SimpleGraph.posSemidef\_lapMatrix}. For $\Lsigned$ we use the
direct sum-of-squares expression
\[
  \pair{x}{\Lsigned x}
  = \!\!\sum_{\{u,v\}\in E\setminus W}\! (x_u - x_v)^2
  + \sum_{\{u,v\}\in W} (x_u + x_v)^2,
\]
which Lemma~\ref{lem:signed-psd} (\S\ref{sec:bounds}) turns into the
Lean fact \texttt{signedLaplacianMobius\_posSemidef}
(\texttt{L13\_PSD.lean}). Equivalently --- but through a heavier route
not used by the Lean proof --- PSD follows from the cover splitting
(Theorem~\ref{thm:cover-splits}): the orientation double cover's
scalar Laplacian is PSD, and similarity to
$\operatorname{diag}(\Lscalar,\Lsigned)$ preserves the PSD property of
each diagonal block.
The Fourier diagonalisation in \S\ref{sec:cycles} below regards both
matrices as acting on $\CC^V$, but the matrices themselves remain
real; complexification is a computational convenience, not a domain
change.

%=============================================================
\section{The Block Decomposition}\label{sec:decomp}
%=============================================================

The connection Laplacian
$L_{\operatorname{Mob}}$ of~\eqref{eq:connection-laplacian}
lives on the product index set $V\times \operatorname{Fin}2$.
Reindexing via the canonical bijection
\[
  \varphi : V\times \operatorname{Fin}2 \;\xrightarrow{\;\sim\;}\; V\sqcup V,
  \qquad (u,0)\mapsto \iota_\ell u,\ (u,1)\mapsto \iota_r u,
\]
we get a matrix on the sum index $V\sqcup V$; the block-diagonal
decomposition then becomes a genuine \emph{block-diagonal} form
\texttt{Matrix.fromBlocks}.

\begin{theorem}\label{thm:decomp}
Let $G$ be a connection graph and $\epsilon \in\{\text{flat},\text{Mob}\}$.
There exists an invertible $P\in\RR^{(V\times\operatorname{Fin}2)
  \times (V\times\operatorname{Fin}2)}$ such that
\[
  \varphi_* \bigl( P\, L_\epsilon\, P^{-1} \bigr)
    = \begin{pmatrix}\Lscalar & 0 \\ 0 & B_\epsilon \end{pmatrix},
    \qquad
  B_\epsilon = \begin{cases}
    \Lsigned & \epsilon=\text{Mob},\\
    \Lscalar & \epsilon=\text{flat.}
  \end{cases}
\]
\end{theorem}

\begin{proof}
Set $P = \II_V \kronprod \widehat R$ (Kronecker product with the
Hadamard~$\widehat R$ of~\eqref{eq:Rhat} acting on the
$\operatorname{Fin}2$ fibre). By functoriality of the Kronecker
product,
\[
  P\cdot P
  = (\II_V\cdot \II_V) \kronprod (\widehat R\,\widehat R)
  = \II_V \kronprod (2\II_2)
  = 2\,\II_{V\times \operatorname{Fin}2},
\]
so $P^{-1} = \tfrac{1}{2}P$, and in particular $\det P\ne 0$.

A direct entrywise calculation using
$\widehat R\,\sigma_x\,\widehat R = 2\diag(1,-1)$
and $\widehat R\,\II_2\,\widehat R = 2\,\II_2$ yields, for
$u,v\in V$ and $i,j\in\operatorname{Fin}2$,
\[
  \tfrac{1}{2}\,\bigl(P\, L_\epsilon\, P\bigr)\bigl[(u,i),(v,j)\bigr]
    = \begin{cases}
      \Lscalar[u,v] & i=j=0,\\
      B_\epsilon[u,v] & i=j=1,\\
      0 & i\ne j.
    \end{cases}
\]
Reindexing via $\varphi$ and using
$P^{-1}=\tfrac{1}{2}P$ turns the left-hand side into
$\varphi_*\bigl(P\,L_\epsilon\,P^{-1}\bigr)$, giving the claim.
The flat case is the special case $W=\emptyset$ of the Mobius
case, or equivalently the fact that $\widehat R\,\II_2\,\widehat R=
2\II_2$ collapses both diagonal blocks to the scalar Laplacian.
\end{proof}

\begin{remark}
In the Lean formalisation \texttt{laplacian\_decomposes} the two
$\epsilon$-branches are proved separately: the Möbius branch uses
$P=\II_V\kronprod \widehat R$ as above, while the flat branch takes
the simpler $P=\II_{V\times\operatorname{Fin}2}$ (identity
similarity) and reduces the claim to
$L_{\mathrm{flat}}=\diag(\Lscalar,\Lscalar)$ after reindex via
$\varphi$. Mathematically the flat case is the specialisation of the
Möbius case at $W=\emptyset$, but the two Lean branches are
syntactically disjoint.
\end{remark}

\begin{remark}
In the Lean script the bijection $\varphi$ is
\texttt{prodFinTwoEquivSum} and the Kronecker-product
identity is the lemma \texttt{mul\_kronecker\_mul} from
\texttt{Mathlib.Data.Matrix.Kronecker}. The entrywise
calculation is made into a single-sum \texttt{Finset.sum\_eq\_single}
reduction that kills $V$-off-diagonal terms using the
$\II_V$ factor of~$P$, and then reduces the remaining
$\operatorname{Fin}2\times \operatorname{Fin}2$ block to one of
three cases by matching on equality of vertices and wrap status.
\end{remark}

%=============================================================
\section{The Orientation Double Cover}\label{sec:cover}
%=============================================================

The preceding decomposition separates the flat part of~$G$ from
its wrap data inside a single graph. The cover perspective lifts
the wrap data into a \emph{branching} structure on a new,
connection-free graph.

\begin{definition}\label{def:cover}
The \emph{orientation double cover} of the connection graph
$(G,W)$ is the simple graph $\tG$ on vertex set
$\tV \defeq V\times \operatorname{Bool}$, with adjacency
\[
  (u,b)\sim_{\tG} (v,c) \iff
    \{u,v\}\in E \;\text{and}\; \bigl(\{u,v\}\in W \Leftrightarrow b\ne c\bigr).
\]
We view $\tG$ as a connection graph with $\widetilde W=\emptyset$:
every edge of $\tG$ is "interior", the cover has \emph{trivial}
connection.
\end{definition}

Equivalently, $\tG$ is the $\ZZ/2$-covering space determined by
the $1$-cochain $\sigma_W:E\to \ZZ/2,\ e\mapsto \mathbf{1}_{e\in W}$
(which is automatically a $1$-cocycle since the graph has no
$2$-cells; see~\S\ref{sec:coho}).
The non-trivial \emph{deck transformation}
$\delta:\tV\to\tV,\ (v,b)\mapsto (v,\neg b)$ is a free graph
involution.

\begin{lemma}\label{lem:deck}
  $\delta$ is a graph automorphism of~$\tG$ of order~$2$ with
  no fixed points.
\end{lemma}

\begin{proof}
Involutivity is immediate from $\neg\neg b = b$. The
adjacency $(\delta x)\sim_{\tG}(\delta y)\Leftrightarrow x\sim_{\tG}y$
follows from the identity
$(\neg b)\ne (\neg c) \Leftrightarrow b\ne c$ and fixity of the
underlying edge $\{u,v\}$. Absence of fixed points is trivial:
$(v,b)\ne (v,\neg b)$. \qed
\end{proof}

\begin{lemma}\label{lem:cover-degree}
  For every $(u,b)\in \tV$, $\deg_{\tG}(u,b) = \deg_G(u)$.
\end{lemma}

\begin{proof}
We construct a bijection $\cN_G(u)\to \cN_{\tG}(u,b)$ by
\[
  v \;\longmapsto\; \bigl(v,\ b\oplus \mathbf{1}_{\{u,v\}\in W}\bigr),
\]
where $\oplus$ is the $\operatorname{Bool}$ XOR. Correctness is the
definition of~$\sim_{\tG}$; injectivity follows by reading the
first coordinate; surjectivity uses the biconditional
$\{u,v\}\in W \Leftrightarrow b\ne c$ to recover the unique~$v$
with $(u,b)\sim_{\tG}(v,c)$.
\end{proof}

We write $\tL\defeq \Lscalar(\tG)$ for the scalar Laplacian of
the cover.

\subsection*{The Hadamard rotation on the $\operatorname{Bool}$-fibre}

Replace the $\operatorname{Fin}2$-fibre of Section~\ref{sec:decomp}
by the $\operatorname{Bool}$-fibre, viewing
\[
  \widehat R : \RR^{\operatorname{Bool}\times \operatorname{Bool}},
  \qquad \widehat R[i,j] = \begin{cases}
    -1 & i=j=\top,\\
    \phantom{-}1 & \text{otherwise},
  \end{cases}
\]
so $\widehat R\,\widehat R = 2\,\II$ on
$\RR^{\operatorname{Bool}\times \operatorname{Bool}}$. Set
\[
  P_{\mathrm{cov}} \defeq \II_V\kronprod \widehat R \quad \in \RR^{\tV\times \tV}.
\]

\begin{theorem}\label{thm:cover-splits}
Under the canonical bijection
$\varphi:\tV\xrightarrow{\sim} V\sqcup V$,
$(v,\bot)\mapsto \iota_\ell v,\ (v,\top)\mapsto \iota_r v$,
we have
\[
  \varphi_*\bigl(P_{\mathrm{cov}}\,\tL\,P_{\mathrm{cov}}^{-1}\bigr)
  = \begin{pmatrix} \Lscalar(G) & 0 \\ 0 & \Lsigned(G)\end{pmatrix}.
\]
In particular, $P_{\mathrm{cov}}$ is invertible and
$\tL\sim \Lscalar(G)\oplus\Lsigned(G)$.
\end{theorem}

\begin{proof}
Write every $2\times 2$ block
\[
  B_{u,v} \defeq \widehat R\,\tL[u,v]\,\widehat R \in
  \RR^{\operatorname{Bool}\times\operatorname{Bool}}
\]
where $\tL[u,v]$ is the $\{u\}\times\operatorname{Bool}$-rows,
$\{v\}\times \operatorname{Bool}$-columns sub-block of $\tL$. A
direct Kronecker calculation gives
\begin{equation}\label{eq:Pcov-identity}
  \bigl(P_{\mathrm{cov}}\,\tL\,P_{\mathrm{cov}}\bigr)\bigl[(u,i),(v,j)\bigr]
    = (\widehat R\,\tL[u,v]\,\widehat R)[i,j]
    = B_{u,v}[i,j].
\end{equation}

Split $B_{u,v}$ into four structural cases:
\begin{enumerate}[label=\textbf{(\alph*)}]
  \item \textit{Diagonal.} If $u=v$, then
  by Lemma~\ref{lem:cover-degree}, $\tL[u,u] = \deg(u)\,\II$,
  whence $\widehat R(\deg(u)\,\II)\widehat R =
  \deg(u)\cdot 2\II$, so $B_{u,u}[i,j] = 2\deg(u)\cdot \delta_{ij}$.
  \item \textit{Adjacent, non-wrap.} If $\{u,v\}\in E\setminus W$,
  then $\tL[u,v]$ has $-1$ on the \emph{diagonal} of
  $\operatorname{Bool}$ (sheets match) and $0$ off-diagonal:
  $\tL[u,v] = (-1)\II$. So
  $B_{u,v}[i,j] = -2\delta_{ij}$.
  \item \textit{Adjacent, wrap.} If $\{u,v\}\in E\cap W$, then
  $\tL[u,v]$ has $-1$ off the diagonal (sheets flip) and $0$ on
  the diagonal: $\tL[u,v] = -\sigma_x$. So
  $\widehat R(-\sigma_x)\widehat R = -2\diag(1,-1)$, i.e.\
  $B_{u,v}[i,j] = -2\delta_{ij}\,(-1)^i$.
  \item \textit{Non-adjacent.} If $\{u,v\}\notin E$, then
  $\tL[u,v]=0$ and $B_{u,v}=0$.
\end{enumerate}

Collecting cases and matching against the defining formulas of
$\Lscalar$ and $\Lsigned$:
\[
  B_{u,v}[i,j] = 2\cdot \delta_{ij}\cdot
  \begin{cases} \Lscalar[u,v] & i=\bot,\\
    \Lsigned[u,v] & i=\top. \end{cases}
\]
Because $P_{\mathrm{cov}}^{-1} = \tfrac{1}{2}P_{\mathrm{cov}}$
(by $P_{\mathrm{cov}}^2 = 2\II$, exactly as in Section~\ref{sec:decomp}),
reindexing through~$\varphi$ gives the claimed block-diagonal
form.
\end{proof}

\begin{corollary}\label{cor:cover-kernel}
$\dim_\RR \Ker\tL = \dim_\RR\Ker \Lscalar(G) +
    \dim_\RR\Ker \Lsigned(G)$.
\end{corollary}

\begin{proof}
Kernel dimension is invariant under similarity by an invertible
matrix and under reindexing by a bijection of the index set; for
a block-diagonal matrix it is additive in the blocks. Apply these
three standard linear-algebra facts to the similarity of
Theorem~\ref{thm:cover-splits}.
\end{proof}

\begin{remark}
In the Lean code these three standard facts are the lemmas
\begin{center}
  \texttt{finrank\_ker\_toLin'\_conj},\quad
  \texttt{finrank\_ker\_toLin'\_reindex},\quad
  \texttt{finrank\_ker\_toLin'\_fromBlocks\_diag}.
\end{center}
They are proved once, over an arbitrary field, in
\texttt{KernelDimension.lean} and reused for both the
block-decomposition Theorem~\ref{thm:decomp} and the
cover-splitting Corollary~\ref{cor:cover-kernel}.
\end{remark}

%=============================================================
\section{$\ZZ/2$ Cohomology and Balanced Components}\label{sec:coho}
%=============================================================

The wrap predicate $W\subseteq E$ reifies canonically as a
$\ZZ/2$-valued $1$-cochain
\[
  \sigma_W : E \;\longrightarrow\; \ZZ/2,\qquad
  \sigma_W(e) \defeq \mathbf{1}_{e\in W}.
\]
It is automatically a $1$-cocycle (there is no $2$-skeleton to
impose constraints on it). Its cohomology class
$[\sigma_W]\in H^1(G;\ZZ/2)$ is the obstruction to
$W$ arising as the coboundary of a vertex $\ZZ/2$-colouring:
$[\sigma_W]=0$ iff there exists $\varepsilon:V\to\operatorname{Bool}$
with $W = \{\,\{u,v\}\in E : \varepsilon(u)\ne \varepsilon(v)\,\}$.
This is the classical notion of \emph{balance} for signed graphs.

\begin{definition}\label{def:balanced}
A connected component $C\in\comp(G)$ is \emph{balanced} iff there
is a map $\varepsilon:V\to\operatorname{Bool}$ such that for every
edge $\{u,v\}\in E$ with $[u]_{\sim_G} = C$,
\[
  \{u,v\}\in W \iff \varepsilon(u)\ne \varepsilon(v).
\]
We write $\comp^{\mathrm{bal}}(G)\subseteq \comp(G)$ for the
subset of balanced components, and
$\beta(G)\defeq \#\comp^{\mathrm{bal}}(G)$.
\end{definition}

\begin{remark}[Scope conventions for Definition~\ref{def:balanced}]
\label{rem:balance-scope}
Two conventions deserve to be flagged.
\begin{enumerate}[label=(\alph*)]
  \item \textbf{Component-local vs.\ graph-global.} The definition
    quantifies only over edges \emph{inside} the component~$C$; the
    values $\varepsilon(w)$ for $w$ outside~$C$ are unconstrained and
    play no role. For a \emph{connected} graph the present definition
    coincides with Harary's original graph-global notion of balance
    ($[\sigma_W]=0$ in $H^1(G;\ZZ/2)$); for disconnected graphs they
    differ, and \emph{every theorem below uses the component-local
    version.} Throughout the paper $\beta(G)$ denotes the
    component-local count.
  \item \textbf{Isolated vertices are vacuously balanced.} If
    $C=\{v_0\}$ is a singleton, the edge quantifier is vacuous and
    the condition holds with any~$\varepsilon$. Consequently every
    isolated vertex counts towards $\beta(G)$; this is load-bearing
    for Theorem~\ref{thm:signed-kernel} below, because an isolated
    vertex contributes a zero row to~$\Lsigned$ and hence~$+1$ to the
    kernel dimension.
\end{enumerate}
\end{remark}

The main combinatorial result of this section relates balance to
the fibres of the induced map on connected components.

\begin{theorem}\label{thm:fiber-card}
Let $\pi:\comp(\tG)\to\comp(G)$ be the map induced by the sheet
projection $\tG\to G,\ (v,b)\mapsto v$. For every
$C\in\comp(G)$,
\[
  \#\pi^{-1}(C) \;=\; \begin{cases} 2 & C \text{ balanced,}\\
    1 & \text{otherwise.} \end{cases}
\]
\end{theorem}

The proof occupies the remainder of this section. Fix
$C\in\comp(G)$ and choose a representative $v_0\in V$ with
$[v_0]_{\sim_G}=C$. Set the \emph{candidate lifts}
\[
  D_\bot \defeq [(v_0,\bot)]_{\comp(\tG)}, \qquad
  D_\top \defeq [(v_0,\top)]_{\comp(\tG)}.
\]

\paragraph{Step 1: every component in the fibre is a candidate lift.}

\begin{lemma}\label{lem:fiber-subset}
For every $D\in\pi^{-1}(C)$, either $D=D_\bot$ or $D=D_\top$.
\end{lemma}

\begin{proof}
Since $\pi([(w,b)]) = [w]_{\comp(G)}$ factors through the first
coordinate, $D=[(w,b)]$ for some $(w,b)\in\tV$ with $[w]=C$. By
connectedness of~$C$ there is a walk $w\to v_0$ in $G$; lifting
edge by edge and setting
\[
  \text{next sheet}\bigl(b,\ \{u,u'\}\bigr)
  \defeq \begin{cases} \neg b & \{u,u'\}\in W,\\
    b & \text{otherwise,} \end{cases}
\]
we obtain a walk $(w,b)\to (v_0,b'')$ in $\tG$ for some
$b''\in\operatorname{Bool}$ determined by the walk. So
$[(w,b)]=[(v_0,b'')]$, and $D\in\{D_\bot,D_\top\}$.
\end{proof}

\paragraph{Step 2: the balanced case.}

\begin{lemma}\label{lem:sep-if-balanced}
If $C$ is balanced, then $D_\bot\ne D_\top$, whence
$\#\pi^{-1}(C)=2$.
\end{lemma}

\begin{proof}
Let $\varepsilon$ witness balance. Define the \emph{XOR}
invariant $\chi:\tV\to \operatorname{Bool}$ by
\[
  \chi(v,b) \defeq \varepsilon(v)\oplus b.
\]
We claim $\chi$ is constant on every connected component
of~$\tG$ whose base image is~$C$. Indeed, for an edge
$(u,b)\sim_{\tG}(u',b')$ with $[u]_{\sim_G}=C$,
\[
  \chi(u,b)\oplus \chi(u',b')
    = \varepsilon(u)\oplus \varepsilon(u')\oplus b\oplus b'
    = \mathbf{1}_{\{u,u'\}\in W} \oplus \mathbf{1}_{\{u,u'\}\in W}
    = 0,
\]
using the definition of~$\sim_{\tG}$ and of balance. So $\chi$
is a locally constant function on the fibre over~$C$, hence
constant on components. But
\[
  \chi(v_0,\bot) = \varepsilon(v_0)\oplus \bot \ne
  \varepsilon(v_0)\oplus \top = \chi(v_0,\top),
\]
so $(v_0,\bot)$ and $(v_0,\top)$ lie in different components:
$D_\bot\ne D_\top$. Together with Lemma~\ref{lem:fiber-subset},
this forces $\pi^{-1}(C) = \{D_\bot,D_\top\}$, of cardinality~$2$.
\end{proof}

\paragraph{Step 3: the unbalanced case.}

\begin{lemma}\label{lem:merge-if-unbalanced}
If $C$ is not balanced, then $D_\bot=D_\top$, whence
$\#\pi^{-1}(C)=1$.
\end{lemma}

\begin{proof}
Unbalance means: for every candidate
$\varepsilon:V\to\operatorname{Bool}$ there is an edge
$\{u,u'\}\in E$ with $[u]=C$ on which
$\mathbf{1}_{\{u,u'\}\in W}\ne \varepsilon(u)\oplus \varepsilon(u')$.
Take the candidate
$\varepsilon_0(v)\defeq \mathbf{1}_{(v,\bot)\sim^*_{\tG}(v_0,\bot)}$
(i.e.\ the characteristic function of the $D_\bot$-fibre
projected to the base). If $D_\bot\ne D_\top$, then $\chi$
constructed from $\varepsilon_0$ as in Lemma~\ref{lem:sep-if-balanced}
would witness balance of~$C$, a contradiction. Hence
$D_\bot = D_\top$, and $\pi^{-1}(C)$ is a singleton.

The Lean formalisation uses the deck transformation~$\delta(v,b)=
(v,\neg b)$ in place of the \emph{ad hoc} $\varepsilon_0$ above: if
$D_\bot=D_\top$ fails then $\delta$ acts freely on $\pi^{-1}(C)$ with
every orbit of size~$2$, forcing $\#\pi^{-1}(C)$ even and the XOR
invariant of Lemma~\ref{lem:sep-if-balanced} becomes a well-defined
$\varepsilon$ witnessing balance of~$C$. Both routes close the same
theorem; the argument is recorded as
\texttt{componentProj\_fiber\_card} in \texttt{L6\_Cohomology.lean}.
\end{proof}

Lemmas~\ref{lem:fiber-subset}--\ref{lem:merge-if-unbalanced} together
prove Theorem~\ref{thm:fiber-card}. An immediate corollary is:

\begin{corollary}\label{cor:components-cover}
$\#\comp(\tG) \;=\; \#\comp(G) \;+\; \beta(G)$.
\end{corollary}

\begin{proof}
Summing $\#\pi^{-1}(C)$ over $C\in\comp(G)$: each balanced
component contributes $2$, each unbalanced contributes~$1$.
So $\#\comp(\tG) = 2\beta(G) + (\#\comp(G)-\beta(G))
  = \#\comp(G)+\beta(G)$.
\end{proof}

%=============================================================
\section{Recognition Theorems}\label{sec:recognition}
%=============================================================

We now reap the kernel-dimension formulas. Throughout, $G$ is a
finite simple connection graph.

\begin{theorem}[Scalar kernel dimension, Mathlib]\label{thm:scalar-ker}
$\dim_\RR \Ker \Lscalar(G) = \#\comp(G)$.
\end{theorem}

This is Mathlib's
\texttt{SimpleGraph.card\_ConnectedComponent\_eq\_rank\_ker\_lapMatrix}
applied to the ordinary graph Laplacian and used as a primitive.

\begin{theorem}\label{thm:signed-kernel}
$\dim_\RR \Ker \Lsigned(G) = \beta(G)$.
\end{theorem}

\begin{proof}
By Corollary~\ref{cor:cover-kernel},
\[
  \dim \Ker\tL \;=\; \dim\Ker\Lscalar(G) + \dim \Ker\Lsigned(G),
\]
where $\tL = \Lscalar(\tG)$. Applying
Theorem~\ref{thm:scalar-ker} to~$\tG$ and to~$G$, and invoking
Corollary~\ref{cor:components-cover},
\[
  \#\comp(G) + \beta(G) = \#\comp(\tG) = \dim\Ker\tL
    = \dim\Ker\Lscalar(G) + \dim\Ker\Lsigned(G)
    = \#\comp(G) + \dim\Ker\Lsigned(G).
\]
Subtracting $\#\comp(G)$ yields the claim.
\end{proof}

\begin{theorem}\label{thm:connection-ker}
For $\epsilon\in\{\mathrm{flat},\mathrm{Mob}\}$,
\[
  \dim_\RR \Ker L_\epsilon
  \;=\; \begin{cases}
    2\,\#\comp(G) & \epsilon=\mathrm{flat},\\
    \#\comp(G)+\beta(G) & \epsilon=\mathrm{Mob}.
  \end{cases}
\]
\end{theorem}

\begin{proof}
By Theorem~\ref{thm:decomp}, $L_\epsilon$ is similar (through the
reindex $\varphi$) to
$\diag(\Lscalar, B_\epsilon)$ with
$B_{\mathrm{flat}}=\Lscalar$, $B_{\mathrm{Mob}}=\Lsigned$. So
\[
  \dim\Ker L_\epsilon
  = \dim\Ker\Lscalar + \dim \Ker B_\epsilon.
\]
Theorems~\ref{thm:scalar-ker} and~\ref{thm:signed-kernel} give the
two diagonal-block kernel dimensions, yielding the claim.
\end{proof}

\begin{theorem}\label{thm:strict}
$\dim\Ker L_{\mathrm{Mob}} < \dim\Ker L_{\mathrm{flat}}$
  iff $\beta(G) < \#\comp(G)$, equivalently, iff \emph{not every}
  component of~$G$ is balanced.
\end{theorem}

\begin{proof}
By Theorem~\ref{thm:connection-ker}, the strict inequality reduces
to $\#\comp(G)+\beta(G) < 2\#\comp(G)$, i.e.\ $\beta(G)<\#\comp(G)$.
The inequality $\beta(G)\le \#\comp(G)$ is immediate from the
inclusion $\comp^{\mathrm{bal}}(G)\subseteq\comp(G)$ of
Definition~\ref{def:balanced}.
\end{proof}

\begin{remark}[Pre-registered plan vs.\ correct statement]
Our pre-registered plan contained the assertion
\[
  \dim\Ker\Lsigned(G) \;\overset{?}{=}\; \#\comp(G) -
    \#\{C\in\comp(G) : |W\cap E(C)|\text{ is odd}\}. \tag{$\star$}
\]
This is \emph{false} in general. Example: let $G=K_4$ with
$W=\{\{1,2\},\{3,4\}\}$, two disjoint wrap edges forming a
matching. The whole graph is one component, the total wrap count
is $2$ (even), so~$(\star)$ predicts
$\dim\Ker\Lsigned = 1$. However a brute-force eigenvalue computation
shows $\dim \Ker\Lsigned = 0$, which is indeed~$\beta(K_4)=0$
from Theorem~\ref{thm:signed-kernel} (the two wrap edges
together form a non-trivial cohomology class on~$K_4$).
Equation~$(\star)$ is correct on $G=C_n$ because
$H^1(C_n;\ZZ/2)\cong\ZZ/2$ is one-dimensional and total
wrap-count mod~$2$ detects the class; it fails as soon as
$\dim H^1(G;\ZZ/2)>1$, i.e.\ as soon as $|E|>|V|-\#\comp(G)$.
\end{remark}

%=============================================================
\section{The Cycle Spectrum}\label{sec:cycles}
%=============================================================

We close with an explicit diagonalisation in the case $G=C_n$, the
$n$-cycle. Identify $V(C_n)=\ZZ/n$, with edges $\{i,i+1\pmod n\}$.
Let $\mathbf{e}_k\in \CC^n$ be the Fourier basis vector
$\mathbf{e}_k[j]= e^{2\pi i k j/n}$ for $k\in\ZZ/n$. Then the
real-valued eigenvectors of any real circulant are
$\operatorname{Re}(\mathbf{e}_k)$ and $\operatorname{Im}(\mathbf{e}_k)$
(with the understanding that $\operatorname{Im}(\mathbf{e}_0)\equiv 0$
and, if $n$ is even, $\operatorname{Im}(\mathbf{e}_{n/2})\equiv 0$,
so those two indices contribute only a single real eigenvector
each).

\begin{theorem}[Flat cycle spectrum]\label{thm:flat-cycle}
The scalar Laplacian $\Lscalar(C_n)$ has spectrum
\[
  \operatorname{Spec}\Lscalar(C_n)
  = \bigl\{\, 2 - 2\cos(2\pi k/n) \,:\, k\in\ZZ/n\,\bigr\},
\]
with $\mathbf{e}_k$ an eigenvector for the eigenvalue
$2-2\cos(2\pi k/n)$.
\end{theorem}

\begin{theorem}[Mobius cycle spectrum]\label{thm:mobius-cycle}
Let $W=\{\{0,1\}\}$ (one wrap edge). Then $\Lsigned(C_n)$ has
spectrum
\[
  \operatorname{Spec}\Lsigned(C_n)
  = \bigl\{\, 2 - 2\cos((2k+1)\pi/n) \,:\, k\in\ZZ/n\,\bigr\}.
\]
\end{theorem}

\begin{proof}[Sketch]
In both cases the Laplacian is a linear combination of the
identity and the shift operator $S:v\mapsto v(\cdot+1)$, so
both commute with the $\ZZ/n$-action and are diagonalised by
the Fourier basis. For $\Lscalar$ the shift weight is $+1$, giving
eigenvalues $2 - (e^{2\pi i k/n}+e^{-2\pi ik/n})$. For $\Lsigned$
with a single wrap edge, a gauge transformation
$v\mapsto \alpha^j v$ with $\alpha=e^{i\pi/n}$ converts the
problem into an identical one with shift weight $-1$, giving
eigenvalues $2 - (e^{(2k+1)\pi i/n}+e^{-(2k+1)\pi i/n})$.
The full Lean proof in \texttt{CycleSpectrum.lean} carries out the
Fourier algebra without appealing to gauge equivalence, producing
each eigenvector explicitly and verifying the eigen-equation by a
direct \texttt{Finset.sum\_eq\_single} reduction.
\end{proof}

\begin{remark}\label{rem:cycle-spectrum-lean}
What \texttt{CycleSpectrum.lean} formalises is the \emph{per-$k$
eigenvalue identity}: for each $k\in\operatorname{Fin} n$ it produces
an explicit nonzero vector $v_k$ with
$\Lscalar(C_n)\cdot v_k=(2-2\cos(2\pi k/n))\,v_k$
(\texttt{scalarCycle\_eigenvalue}, $n\ge 3$) and similarly for the
signed Mobius cycle (\texttt{signedCycle\_eigenvalue}). Assembly of
these per-$k$ identities into the full multiset equality
$\operatorname{Spec}\Lscalar(C_n)=\{2-2\cos(2\pi k/n)\}$ of
Theorems~\ref{thm:flat-cycle}--\ref{thm:mobius-cycle} follows by a
dimension count, and is not itself a named Lean theorem.
\end{remark}

\begin{corollary}
On $G=C_n$, the Mobius kernel strictly drops iff $n$ is odd (so that
$\{2-2\cos((2k+1)\pi/n)\}$ is positive for all~$k$), matching
$\beta(C_n)=0$ when $W=\{\{0,1\}\}$ and~$n$ is odd.
\end{corollary}

%=============================================================
\section{The Lean Formalisation}\label{sec:formalisation}
%=============================================================

The entire development is formalised in Lean~4 against
Mathlib~v4.11.0. The dependency graph is linear:
\[
  \text{\texttt{Basic}} \to
  \text{\texttt{KernelDimension}} \to
  \text{\texttt{L5\_Cover}} \to
  \text{\texttt{L6\_Cohomology}} \to
  \text{\texttt{L8\_Recognition}},
\]
with a lateral \texttt{CycleSpectrum} file depending on
\texttt{Basic} and \texttt{KernelDimension}.

\subsection*{Axiomatic footprint}

Every named theorem of this paper was checked by the
\texttt{\#print axioms} command to depend only on the three
standard Lean axioms \texttt{propext}, \texttt{Classical.choice},
and \texttt{Quot.sound}; none involve \texttt{sorryAx} or any
user-introduced axiom. In particular:
\begin{itemize}
  \item Theorem~\ref{thm:decomp} is
    \texttt{ConnGraph.laplacian\_decomposes};
  \item Theorem~\ref{thm:cover-splits} is
    \texttt{ConnGraph.scalarLap\_cover\_splits};
  \item Theorem~\ref{thm:fiber-card} is
    \texttt{ConnGraph.componentProj\_fiber\_card};
  \item Theorem~\ref{thm:signed-kernel} is
    \texttt{ConnGraph.signedLaplacian\_kernel\_dim\_general};
  \item Theorem~\ref{thm:connection-ker} is
    \texttt{ConnGraph.connectionLaplacian\_kernel\_dim\_general};
  \item Theorem~\ref{thm:strict} is
    \texttt{ConnGraph.mobius\_kernel\_strict\_iff\_general};
  \item Theorems~\ref{thm:flat-cycle}--\ref{thm:mobius-cycle} are
    certified per eigenvalue by
    \texttt{CycleSpectrum.scalarCycle\_eigenvalue} and
    \texttt{CycleSpectrum.signedCycle\_eigenvalue}; see
    Remark~\ref{rem:cycle-spectrum-lean} below.
\end{itemize}

\subsection*{Numerical cross-check}

As an independent safety net against the formal statements being
merely \emph{provable} but not intended, the accompanying Python
script \texttt{findings/round2/fuzzer/fuzz.py} enumerates every
simple graph on
$n\in\{2,3,4,5\}$ vertices up to isomorphism (51~graphs, matching
OEIS A000088) and every $2^{|E|}$ wrap subset, totalling
$3{,}456$~configurations. For each~$(G,W)$ it checks:
\begin{itemize}
  \item Theorem~\ref{thm:cover-splits} via spectrum equality of
    $\tL$ and $\Lscalar(G)\oplus \Lsigned(G)$;
  \item Corollary~\ref{cor:cover-kernel} via SVD-based integer
    kernel-dimension arithmetic;
  \item Theorem~\ref{thm:fiber-card} via brute-force
    $\ZZ/2$-colouring search and connected-component counting.
\end{itemize}
All $3{,}456$ configurations agreed with the formal theorem
statements; zero disagreements were found.

\subsection*{Summary of effort distribution}

The ``mathematical'' content, in descending order of novelty:
\begin{enumerate}[label=\textbf{(\arabic*)}]
  \item The fibre-cardinality theorem
    (Theorem~\ref{thm:fiber-card}) has no counterpart in
    Mathlib; it requires the walk-lift construction,
    XOR-parity invariant, and deck-action freeness argument,
    all of which are built from scratch against Mathlib's
    \texttt{SimpleGraph.ConnectedComponent} eliminators.
  \item The cover-splitting Theorem~\ref{thm:cover-splits} is an
    entrywise matrix identity. Mathlib's Kronecker library
    (\texttt{Matrix.mul\_kronecker\_mul}, etc.) mechanises
    $\sim\!15\%$; the remaining work is a four-case matching on
    wrap/non-wrap adjacency and a 2x2 block calculation.
  \item The block-decomposition Theorem~\ref{thm:decomp} reuses
    the Kronecker pattern from (2) on the $\operatorname{Fin}2$
    fibre instead of the $\operatorname{Bool}$ fibre.
  \item Recognition Theorems~\ref{thm:signed-kernel} and
    \ref{thm:connection-ker} are a few lines of arithmetic
    conclusion once (1)--(3) are in place.
\end{enumerate}

\subsection*{Lessons}

The formalisation surfaced one genuine bug in our informal plan
(the false equation~$(\star)$). This is a modest-but-real example
of Lean's mathematical utility: the computer cannot be negotiated
with when $\beta(G)$ and
$\#\comp(G)-\#\{\text{odd-wrap components}\}$ fail to agree on
$K_4$ with a matching wrap.

\bigskip

\noindent\textit{Acknowledgements.} This work was carried out as
an independent formal-verification project on combinatorial
spectral identities, under a pre-registered discovery and
adversarial-review discipline.

%=============================================================
\section{Further Identities: Traces, Bounds, and the
Characteristic Polynomial}\label{sec:bounds}
%=============================================================

The block decomposition (Theorem~\ref{thm:decomp}) and cover splitting
(Theorem~\ref{thm:cover-splits}) imply a family of exact spectral
identities that sit one inference-step away from the recognition
theorems. We collect the ones that admit a one-paragraph proof and
are stated at the level of generality on which the formalised
library already operates. Together they close the
``trace-and-inequality gap'' left open in earlier versions of this
manuscript. All are formalised in
\texttt{ConnectionLaplacian/L9\_Bounds.lean}, and the
multiplicative spectral partition
(Theorem~\ref{thm:cover-charpoly}) is formalised in
\texttt{ConnectionLaplacian/L10\_CoverCharpoly.lean}.

\begin{proposition}[Trace formulas]\label{prop:traces}
For every connection graph $G=(V,E,W)$,
\begin{align*}
  \tr \Lscalar(G) \;&=\; 2\,\#E, &
  \tr \Lsigned(G) \;&=\; 2\,\#E, \\
  \tr L_{\mathrm{flat}}(G) \;&=\; 4\,\#E, &
  \tr L_{\mathrm{Mob}}(G) \;&=\; 4\,\#E.
\end{align*}
\end{proposition}

\begin{proof}
The diagonal entry of $\Lscalar$ at $v$ is
$\deg(v)$; irreflexivity of the underlying simple graph kills the
wrap off-diagonal contribution at the diagonal, so the diagonal of
$\Lsigned$ is also $\deg(v)$; and the diagonals of
$L_{\mathrm{flat}}, L_{\mathrm{Mob}}$ are each $\deg(v)\cdot\II_2$,
contributing $2\deg(v)$. Summing over $v$ and applying the
degree-sum identity
$\sum_v \deg(v) = 2\,\#E$
(Mathlib's
\texttt{SimpleGraph.sum\_degrees\_eq\_twice\_card\_edges}) proves
all four equalities.
\end{proof}

\begin{proposition}[Quadratic trace and Frobenius
identities]\label{prop:traces2}
For every connection graph $G=(V,E,W)$,
\begin{align*}
  \tr\bigl(L_{\mathrm{Mob}}(G)^2\bigr)
  \;&=\; \tr\bigl(\Lscalar(G)^2\bigr)
       + \tr\bigl(\Lsigned(G)^2\bigr), \\
  \|L_{\mathrm{Mob}}(G)\|_F^2
  \;&=\; \|\Lscalar(G)\|_F^2 + \|\Lsigned(G)\|_F^2,
\end{align*}
and the cross term obeys the polarisation identity
\[
  2\,\tr\!\bigl(\Lscalar(G)\cdot\Lsigned(G)\bigr)
  \;=\;
  \tr\!\bigl((\Lscalar(G){+}\Lsigned(G))^2\bigr)
  - \tr\!\bigl(\Lscalar(G)^2\bigr)
  - \tr\!\bigl(\Lsigned(G)^2\bigr).
\]
\end{proposition}

\begin{proof}
Starting from the Kronecker-conjugation similarity of
Theorem~\ref{thm:decomp} in its Möbius form, the trace-of-square and
Frobenius identities follow from invariance of $\tr(M^2)$ and of
$\tr(M^\top M)$ under a same-sided reindex composed with the
involution $P=\II_V\kronprod\widehat R$ with $P^{-1}=\tfrac12 P$.
The cross-term identity is pure algebra: expand
$(\Lscalar+\Lsigned)^2$ and use symmetry of the scalar and signed
Laplacians together with the cyclic property
$\tr(\Lscalar\Lsigned)=\tr(\Lsigned\Lscalar)$. The Lean names are
\texttt{trace\_sq\_laplacian\_decomposes},
\texttt{frobenius\_laplacian\_decomposes} and
\texttt{trace\_mul\_scalar\_signed\_eq}.
\end{proof}

\begin{remark}[Polarisation versus combinatorial form]
The cross-term identity above is stated in \emph{polarisation form}
and is pure matrix algebra. An equivalent combinatorial form
$\tr(\Lscalar\cdot\Lsigned)=\tr({\Lscalar}^{2})-4\,\#W$, which can be
read off the entrywise formula for $\Lscalar\cdot\Lsigned$ by
summing contributions at diagonal and off-diagonal positions, is
confirmed numerically by the R6 Stage-6 fuzzer (zero deviation
across $300$ random connected graphs at $n=7$) but is not yet a
separately named Lean theorem; a $\le 20$-line bridge lemma is
recorded as an R7 carry-over.
\end{remark}

\begin{proposition}[Kernel inequalities]\label{prop:kernel-ineq}
For every connection graph $G$,
\begin{align*}
  \dim\Ker \Lsigned(G) \;&\le\; \dim\Ker \Lscalar(G), \\
  \dim\Ker L_{\mathrm{flat}}(G) - \dim\Ker L_{\mathrm{Mob}}(G)
    \;&=\; \#\comp(G) - \beta(G),
\end{align*}
and in particular the difference is bounded by $\#\comp(G)$.
\end{proposition}

\begin{proof}
By Theorems~\ref{thm:scalar-ker}, \ref{thm:signed-kernel},
\ref{thm:connection-ker}, the LHS of the first line equals
$\beta(G)$, the RHS equals $\#\comp(G)$, and
$\beta(G)\le \#\comp(G)$ because
$\comp^{\mathrm{bal}}(G)\subseteq \comp(G)$. For the second line,
$\dim\Ker L_{\mathrm{flat}} = 2\,\#\comp(G)$ and
$\dim\Ker L_{\mathrm{Mob}} = \#\comp(G)+\beta(G)$, whose difference
is $\#\comp(G)-\beta(G)$, itself bounded by $\#\comp(G)$.
\end{proof}

\begin{lemma}[Positive semidefiniteness of the signed
Laplacian]\label{lem:signed-psd}
$\Lsigned(G) \succeq 0$. Indeed
\[
  \pair{x}{\Lsigned(G)\, x}
  = \!\!\sum_{e=\{u,v\}\in E\setminus W}\! (x_u - x_v)^2
  + \sum_{e=\{u,v\}\in W} (x_u + x_v)^2
\qquad \text{for all } x\in\RR^V.
\]
\end{lemma}

\begin{proof}
Fix $x\in\RR^V$. The off-diagonal entry of $\Lsigned$ at $(u,v)$
equals $+1$ when $\{u,v\}\in W$ and $-1$ when $\{u,v\}\in E\setminus W$;
the diagonal is $\deg(u)$. Since $\Lsigned$ is symmetric, each
unordered edge $\{u,v\}$ contributes the pair of entries $(u,v)$
and $(v,u)$ to the quadratic form, yielding a factor of~$2$:
\[
  \pair{x}{\Lsigned x}
  \;=\; \sum_u \deg(u)\,x_u^2
      \;+\; 2\!\!\sum_{\{u,v\}\in W}\!\!x_u x_v
      \;-\; 2\!\!\sum_{\{u,v\}\in E\setminus W}\!\!x_u x_v.
\]
Rewriting $\sum_u \deg(u)\,x_u^2
= \sum_{\{u,v\}\in E}\bigl(x_u^2 + x_v^2\bigr)$ (each edge
contributes to the degree of both endpoints) and splitting the
degree sum along $E = (E\setminus W)\sqcup W$:
\[
  \pair{x}{\Lsigned x}
  \;=\;\!\!\sum_{\{u,v\}\in E\setminus W}\!\!(x_u^2 + x_v^2 - 2 x_u x_v)
    \;+\;\sum_{\{u,v\}\in W}(x_u^2 + x_v^2 + 2 x_u x_v),
\]
which is exactly the announced sum of squares. Every summand is
$\ge 0$, so $\pair{x}{\Lsigned x}\ge 0$.
\end{proof}

\begin{theorem}[Multiplicative spectral partition: cover
characteristic polynomial]\label{thm:cover-charpoly}
Let $\tL = \Lscalar(\tG)$ be the scalar Laplacian of the orientation
double cover of~$G$. Then the characteristic polynomials factor as
\[
  \det\bigl(\lambda\,\II_{\tV} - \tL\bigr) \;=\;
  \det\bigl(\lambda\,\II_V - \Lscalar(G)\bigr)\cdot
  \det\bigl(\lambda\,\II_V - \Lsigned(G)\bigr).
\]
Equivalently, as multisets the spectra satisfy
$\operatorname{Spec}(\tL) = \operatorname{Spec}(\Lscalar(G)) \uplus
\operatorname{Spec}(\Lsigned(G))$.
\end{theorem}

\begin{remark}[Direct-bundle factorisation]\label{rem:direct-bundle-charpoly}
Composing Theorem~\ref{thm:cover-charpoly} with the similarity of
Theorem~\ref{thm:decomp} gives the same factorisation directly for
$L_{\mathrm{Mob}}$ without invoking the cover:
\[
  \det\bigl(\lambda\,\II - L_{\mathrm{Mob}}(G)\bigr)
  \;=\;
  \det\bigl(\lambda\,\II - \Lscalar(G)\bigr)\cdot
  \det\bigl(\lambda\,\II - \Lsigned(G)\bigr),
\]
and for every $\alpha\in\RR$,
$\det(L_{\mathrm{Mob}}+\alpha\II)=\det(\Lscalar+\alpha\II)\cdot
\det(\Lsigned+\alpha\II)$. As multisets,
$\operatorname{Spec}(L_{\mathrm{Mob}}) = \operatorname{Spec}(\Lscalar)
\uplus \operatorname{Spec}(\Lsigned)$. These are Lean theorems
\texttt{mobius\_charpoly\_eq\_scalar\_times\_signed},
\texttt{shiftedDet\_factorises} and
\texttt{mobius\_charpoly\_roots\_eq\_union}.
\end{remark}

\begin{proof}
By Theorem~\ref{thm:cover-splits}, $\tL$ is similar (through an
invertible Hadamard--Kronecker change of basis composed with the
reindex $V\times\operatorname{Bool}\cong V\sqcup V$) to the block
matrix $\diag(\Lscalar(G),\Lsigned(G))$. Similarity preserves the
characteristic polynomial, and the characteristic polynomial of a
block-diagonal matrix factors into the product of its diagonal
blocks' characteristic polynomials. The multiset equality follows by
equating roots with multiplicity.
\end{proof}

\begin{remark}[Interlacing as a corollary]
Theorem~\ref{thm:cover-charpoly} strictly strengthens the Cauchy
interlacing one might expect between the spectra of $\Lsigned$ and
$\tL$: interlacing would merely bound each eigenvalue of $\Lsigned$
between two eigenvalues of $\tL$, whereas
Theorem~\ref{thm:cover-charpoly} pins down the spectrum of $\tL$
exactly as the disjoint union of the two smaller spectra. In
particular $\det\bigl(\lambda\II - L_{\mathrm{Mob}}(G)\bigr) =
\det\bigl(\lambda\II - L_{\mathrm{flat}}(G)\bigr)$ whenever every
component of~$G$ is balanced.
\end{remark}

\begin{corollary}[Pointwise multiplicity]\label{cor:mult-pointwise}
For every $\lambda\in\RR$,
\[
  m_\lambda(\tL) \;=\; m_\lambda(\Lscalar(G)) + m_\lambda(\Lsigned(G)),
\]
where $m_\lambda(M)$ denotes the multiplicity of~$\lambda$ as an
eigenvalue of~$M$.
\end{corollary}

\begin{proof}
Immediate from Theorem~\ref{thm:cover-charpoly}: the multiplicity of
$\lambda$ in a product of polynomials is the sum of multiplicities
in the factors.
\end{proof}

%=============================================================
\section{Structure Theorems for Trees, Forests, and Cycles}
\label{sec:structure}
%=============================================================

We record the ``topological'' side of the recognition theorems: on
graphs where $H^1$ is understood, the content of
Theorem~\ref{thm:signed-kernel} becomes explicit.

\begin{theorem}[Tree balance]\label{thm:tree-balanced}
If the underlying simple graph of~$G$ is a tree, then
every component of~$G$ is balanced, for \emph{every} choice of wrap
set~$W\subseteq E(G)$. Consequently
\[
  \beta(G) = \#\comp(G), \qquad
  \dim\Ker \Lsigned(G) = \#\comp(G), \qquad
  \dim\Ker L_{\mathrm{Mob}}(G) = 2\,\#\comp(G),
\]
and the strict inequality of Theorem~\ref{thm:strict} fails.
\end{theorem}

\begin{proof}
Fix a root $r$ of each tree component. For $v$ in that component,
let $p_{r\to v}$ be the unique simple $r\to v$ path (Mathlib's
\texttt{SimpleGraph.IsAcyclic.path\_unique}). Define
$\varepsilon(v)\defeq \sum_{e\in p_{r\to v}} \mathbf{1}_{e\in W}
  \pmod 2$, with $\varepsilon(r)=\bot$ by the empty sum. We show
$\varepsilon(u)\oplus \varepsilon(v) = \mathbf{1}_{\{u,v\}\in W}$
for every edge $\{u,v\}\in E$ in the component.

Consider the closed walk $w \defeq p_{r\to u} \cdot \{u,v\} \cdot
p_{v\to r}$. Because the underlying graph is acyclic, the only
closed walks at~$r$ are back-tracking compositions of simple edges,
so $\sum_{e\in w}\mathbf{1}_{e\in W}\equiv 0\pmod 2$ (every
traversed edge is traversed an even number of times). Splitting
the sum into the $p_{r\to u}$-part, the single edge $\{u,v\}$, and
the $p_{v\to r}$-part (the latter summing to $\varepsilon(v)$ by
symmetry of $\mathbf{1}_{\cdot\in W}$) gives
\[
  0 \;\equiv\; \varepsilon(u) + \mathbf{1}_{\{u,v\}\in W}
  + \varepsilon(v) \pmod 2,
\]
hence $\varepsilon(u)\oplus\varepsilon(v)=\mathbf{1}_{\{u,v\}\in W}$.
Thus $\varepsilon$ witnesses balance
(Definition~\ref{def:balanced}), and every tree component is
balanced. The kernel-dimension consequences are immediate from
Theorems~\ref{thm:signed-kernel} and~\ref{thm:connection-ker}.
\end{proof}

\begin{corollary}\label{cor:forest-star-path}
If the underlying graph is a connected tree, then
$\dim\Ker L_{\mathrm{Mob}}(G) = 2$ for every~$W$. In
particular paths $P_n$ and stars $K_{1,n}$ have kernel dimension
$2$ independently of~$W$. The statement generalises to disconnected
forests with $\dim\Ker L_{\mathrm{Mob}}(G)=2\,\#\comp(G)$ by applying
the connected case per-component, but this lift is not presently a
separately named Lean theorem --- only the \texttt{IsTree}
(connected-acyclic) case is formalised.
\end{corollary}

\begin{theorem}[Walk-sum detector for balance]\label{thm:walk-sum}
Let $C\in\comp(G)$. Then $C$ is balanced iff every closed walk
$w:v\to v$ with $v\in C$ satisfies
$\sum_{e\in w} \mathbf{1}_{e\in W} \equiv 0 \pmod 2$.
\end{theorem}

\begin{proof}
$(\Rightarrow)$ If $\varepsilon$ witnesses balance then for each
edge $e=\{u,u'\}$ in~$C$, $\varepsilon(u)\oplus \varepsilon(u') =
\mathbf{1}_{e\in W}$. Summing along a closed walk $v\to v$ the
$\varepsilon$-telescopes, leaving
$\sum_e \mathbf{1}_{e\in W}\equiv 0$.
$(\Leftarrow)$ Pick $v_0\in C$ and define
$\varepsilon(v)\defeq \sum_{e\in w} \mathbf{1}_{e\in W}\pmod 2$
for any walk $w:v_0\to v$; the closed-walk-even hypothesis makes
$\varepsilon$ walk-independent. By construction
$\varepsilon(u)\oplus \varepsilon(v)\equiv \mathbf{1}_{e\in W}$ on
every edge $e=\{u,v\}$ of~$C$, so $\varepsilon$ witnesses balance.
\end{proof}

\begin{corollary}[Cycle even-wrap formula]\label{cor:cycle-ew}
If $G=C_n$ is a single cycle of length~$n\ge 3$, then
\[
  \beta(C_n,W) = \begin{cases} 1 & \#W \equiv 0\pmod 2,\\
    0 & \#W \equiv 1\pmod 2,\end{cases}
  \qquad
  \dim\Ker \Lsigned(C_n,W) = \begin{cases} 1 & \#W \text{ even,}\\
    0 & \#W \text{ odd.}\end{cases}
\]
In particular the pre-registered formula
$\dim\Ker\Lsigned = \#\comp(G) - \#\{\text{odd-wrap components}\}$
holds exactly for cycle graphs.
\end{corollary}

\begin{proof}
The first homology of an $n$-cycle with $\ZZ/2$-coefficients is
one-dimensional, $H_1(C_n;\ZZ/2)\cong\ZZ/2$, generated by the
fundamental cycle~$\gamma$. Every closed walk $w$ at a basepoint
decomposes, modulo back-tracking (which contributes
edges twice and hence~$0$ to the $\ZZ/2$-sum of wrap indicators),
into an integer multiple of~$\gamma$; its wrap-parity therefore
equals $[w:\gamma]\cdot\#W \pmod 2$. The walk-sum detector of
Theorem~\ref{thm:walk-sum} therefore reduces to the parity
of~$\#W$ and Theorem~\ref{thm:signed-kernel} converts
$\beta\in\{0,1\}$ to $\dim\Ker\Lsigned\in\{0,1\}$.
\end{proof}

\begin{remark}
The Lean certificate
\texttt{cycle\_isBalanced\_iff\_even\_wrap\_weak}
(\texttt{L14\_CycleEw.lean}) names this a \emph{weak} form: it
takes an Eulerian closed walk around the cycle and a reverse-direction
hypothesis as explicit inputs to the caller. Lifting this to the
fully constructive form stated above requires a
\texttt{fundamentalCycleWalk} helper plus a winding-number argument
currently tracked as an R7 carry-over.
\end{remark}

%=============================================================
\section{Counter-claims: What the Formalisation Does
\emph{Not} Prove}\label{sec:counter}
%=============================================================

The extended fuzzer (\S\ref{sec:formalisation}) and the
\texttt{NEGATOR} pass of our discovery swarm refuted four naive
statements that were, at various points, attempted drafts of the
identities above. We record them here both for historical honesty
and as explicit boundary markers for the theorems that \emph{did}
survive.

\begin{example}[Naive equality of M\"obius and scalar kernels is
false]\label{ex:naive-eq}
A tempting but wrong reformulation of Theorem~\ref{thm:strict} is:
``$\dim\Ker L_{\mathrm{Mob}} = \dim\Ker L_{\mathrm{scalar}}$ iff
every component of~$G$ is balanced.'' This fails already on the
empty graph on $n=3$ vertices with $W=\emptyset$: every
component is trivially balanced (vacuously so, cf.\
Remark~\ref{rem:balance-scope}), yet
$\dim\Ker L_{\mathrm{Mob}} = 6$ (a kernel inside $\CC^{2V}$) while
$\dim\Ker L_{\mathrm{scalar}} = 3$ (a kernel inside $\CC^{V}$): the
two kernels sit in different ambient spaces, so the attempted
identity already fails on dimensional grounds. The correct
equivalence uses
the \emph{flat} Laplacian: $\dim\Ker L_{\mathrm{Mob}} =
\dim\Ker L_{\mathrm{flat}} = 2\,\#\comp(G)$ iff every component is
balanced, which is Theorem~\ref{thm:connection-ker} combined with
Theorem~\ref{thm:scalar-ker}.
\end{example}

\begin{example}[A single wrap edge can leave the graph
balanced]\label{ex:bridge-wrap}
The implication ``$G$ connected and $\#W=1$ $\Rightarrow$ $G$ is not
balanced'' is false. Counterexample: the path $P_3$ with vertices
$v_0\text{-}v_1\text{-}v_2$ and the single wrap edge
$\{v_0,v_1\}$ (a bridge) is balanced: the assignment
$\varepsilon(v_0)=\top,\ \varepsilon(v_1)=\bot,\
\varepsilon(v_2)=\bot$ witnesses balance, since
$\varepsilon(v_0)\oplus\varepsilon(v_1)=1=\mathbf{1}_{\{v_0,v_1\}\in W}$
and $\varepsilon(v_1)\oplus\varepsilon(v_2)=0=
\mathbf{1}_{\{v_1,v_2\}\in W}$. The correct form is: a single wrap
edge unbalances the graph iff the wrap edge lies on a cycle.
\end{example}

\begin{example}[Adding a non-wrap edge between components can
\emph{decrease} $\beta$]\label{ex:cross-merge}
The fuzzy-logic sweep (\S\ref{sec:formalisation}, claim
\texttt{C2}, $\tau=0.0000$) refutes the naive monotonicity
``adding a non-wrap edge can only preserve or increase
$\beta$''. On three isolated vertices ($W=\emptyset$, so
$\beta=3$), adding the non-wrap edge $\{v_0,v_1\}$ merges two
balanced components into one still-balanced component and drops
$\beta$ to $2$. The correct law is bidirectional: for a non-wrap
edge $e$ joining two distinct components,
$\beta(G+e)\in\{\beta(G)-1,\beta(G)\}$, with $-1$ when \emph{at
least one} endpoint lies in a balanced component and $0$ when both
endpoints lie in unbalanced components. (In the mixed
balanced-plus-unbalanced case, the merged component is unbalanced,
so the pre-existing balanced component is lost and $\beta$ drops
by~$1$.)
\end{example}

\begin{example}[Non-wrap edge contraction is not kernel-stable]
\label{ex:contraction}
On $K_3$ with $W=\{\{v_0,v_1\}\}$ the signed Laplacian has
kernel dimension~$0$. Contracting the non-wrap edge $\{v_0,v_2\}$
identifies $v_0$ and $v_2$, leaving the single wrap edge
$\{v_0,v_1\}$, a tree on $2$ vertices: the signed Laplacian
now has kernel dimension~$1$. So non-wrap edge contraction does
\emph{not} preserve $\dim\Ker\Lsigned$ in general.
\end{example}

%=============================================================
\appendix
\section{Formalisation index}\label{app:formal-index}
%=============================================================

The principal named theorems, propositions, corollaries, and
lemmas of the main text are paired below with corresponding
Lean~4 declarations in the
\texttt{connection\_laplacian\_lean} sandbox. (A handful of
subsidiary lemmas used only in intermediate proofs---e.g.\ the
deck-transformation Lemma~\ref{lem:deck}, the cover-degree
Lemma~\ref{lem:cover-degree}, and the fiber-subset/separation
Lemmas~\ref{lem:fiber-subset}, \ref{lem:sep-if-balanced},
\ref{lem:merge-if-unbalanced}---are implicitly witnessed inside
the proofs of theorems listed below and are not tabulated
separately.) The table below
pairs them explicitly. All listed declarations are sorry-free and
their \texttt{\#print axioms} output is exactly
$\{\texttt{propext},\, \texttt{Classical.choice},\, \texttt{Quot.sound}\}$ —
the standard constructive--classical foundation of Mathlib, with no
\texttt{sorryAx} and no ad-hoc \texttt{axiom}~declarations. The file
paths are relative to
\texttt{connection\_laplacian\_lean/ConnectionLaplacian/}.

\begin{center}\footnotesize
\renewcommand{\arraystretch}{1.05}
\begin{tabular}{@{}p{3.2cm}p{6.2cm}p{4.4cm}@{}}
\toprule
Paper statement & Lean name & File \\
\midrule
Thm.~\ref{thm:scalar-ker} &
  \texttt{scalarLaplacian\_kernel\_dim} & \texttt{KernelDim.lean} \\
Thm.~\ref{thm:decomp} (direct sum) &
  \texttt{laplacian\_decomposes} & \texttt{KernelDim.lean} \\
Thm.~\ref{thm:decomp} (ker dim) &
  \texttt{laplacian\_kernel\_dim\_decomposes} & \texttt{L8\_Recognition.lean} \\
\S\ref{sec:cover} &
  \texttt{scalarLap\_cover\_splits} & \texttt{L5\_Cover.lean} \\
Def.~\ref{def:balanced} &
  \texttt{isBalanced} & \texttt{L6\_Cohomology.lean} \\
Thm.~\ref{thm:fiber-card} &
  \texttt{componentProj\_fiber\_card} & \texttt{L6\_Cohomology.lean} \\
Thm.~\ref{thm:signed-kernel} &
  \texttt{signedLaplacian\_kernel\_dim\_general} & \texttt{L8\_Recognition.lean} \\
Thm.~\ref{thm:connection-ker} &
  \texttt{connectionLaplacian\_kernel\_dim\_general} & \texttt{L8\_Recognition.lean} \\
Thm.~\ref{thm:strict} &
  \texttt{mobius\_kernel\_strict\_iff\_general} & \texttt{L8\_Recognition.lean} \\
Prop.~\ref{prop:traces} (sc) &
  \texttt{trace\_scalarLaplacian} & \texttt{L9\_Bounds.lean} \\
Prop.~\ref{prop:traces} (sg) &
  \texttt{trace\_signedLaplacianMobius} & \texttt{L9\_Bounds.lean} \\
Prop.~\ref{prop:kernel-ineq} (mono) &
  \texttt{signed\_kernel\_le\_scalar\_kernel} & \texttt{L9\_Bounds.lean} \\
Prop.~\ref{prop:kernel-ineq} (drop~$=\beta$) &
  \texttt{kernel\_drop\_eq\_unbalanced} & \texttt{L9\_Bounds.lean} \\
Prop.~\ref{prop:kernel-ineq} (drop~$\le\#\comp$) &
  \texttt{kernel\_drop\_le\_numComponents} & \texttt{L9\_Bounds.lean} \\
Lem.~\ref{lem:signed-psd} &
  \texttt{signedLaplacianMobius\_posSemidef} & \texttt{L13\_PSD.lean} \\
Thm.~\ref{thm:cover-charpoly} &
  \texttt{cover\_charpoly\_eq\_scalar\_times\_mobius} & \texttt{L10\_CoverCharpoly.lean} \\
Cor.~\ref{cor:mult-pointwise} &
  \multicolumn{2}{l}{\emph{corollary of \texttt{cover\_charpoly\_eq\_scalar\_times\_mobius}}} \\
Thm.~\ref{thm:tree-balanced} &
  \texttt{tree\_isBalanced\_of\_isTree} & \texttt{L11\_Trees.lean} \\
\quad (ker dim) &
  \texttt{connectionLaplacian\_kernel\_dim\_tree\_mobius} & \texttt{L11\_Trees.lean} \\
\quad (gap zero) &
  \texttt{tree\_kernel\_gap\_zero} & \texttt{L11\_Trees.lean} \\
Cor.~\ref{cor:forest-star-path} &
  \multicolumn{2}{l}{\emph{specialisation of the tree ker-dim theorem}} \\
Thm.~\ref{thm:walk-sum} &
  \texttt{isBalanced\_iff\_closedWalk\_wrap\_even} & \texttt{L12\_WalkH1.lean} \\
Cor.~\ref{cor:cycle-ew} &
  \texttt{cycle\_isBalanced\_iff\_even\_wrap\_weak} & \texttt{L14\_CycleEw.lean} \\
Ex.~\ref{ex:cross-merge} (R5 refinement, C1 $\tau=1$) &
  \texttt{numBalancedComponents\_monotone\_remove\_nonwrap\_nonbridge} & \texttt{L15\_BridgeMonotone.lean} \\
Rem.~\ref{rem:direct-bundle-charpoly} (direct charpoly) &
  \texttt{mobius\_charpoly\_eq\_scalar\_times\_signed} & \texttt{L16\_SpectrumUnion.lean} \\
Rem.~\ref{rem:direct-bundle-charpoly} (spec union) &
  \texttt{mobius\_charpoly\_roots\_eq\_union} & \texttt{L16\_SpectrumUnion.lean} \\
Rem.~\ref{rem:direct-bundle-charpoly} (shifted det) &
  \texttt{shiftedDet\_factorises} & \texttt{L17\_TracesAndLipschitz.lean} \\
Prop.~\ref{prop:traces} ($\tr L_{\mathrm{Mob}}$) &
  \texttt{trace\_laplacian\_mobius} & \texttt{L17\_TracesAndLipschitz.lean} \\
Prop.~\ref{prop:traces2} ($\tr(\cdot^2)$) &
  \texttt{trace\_sq\_laplacian\_decomposes} & \texttt{L17\_TracesAndLipschitz.lean} \\
Prop.~\ref{prop:traces2} (Frobenius) &
  \texttt{frobenius\_laplacian\_decomposes} & \texttt{L17\_TracesAndLipschitz.lean} \\
Prop.~\ref{prop:traces2} (polarisation) &
  \texttt{trace\_mul\_scalar\_signed\_eq} & \texttt{L17\_TracesAndLipschitz.lean} \\
\bottomrule
\end{tabular}
\end{center}
\noindent\footnotesize (File paths: \texttt{KernelDim.lean} abbreviates \texttt{KernelDimension.lean}.)\normalsize

\paragraph{Reproducibility.} From the project root,
\begin{quote}
\verb|elan run lake build|
\end{quote}
rebuilds all 1833 Lean objects (Mathlib~v4.11.0 +
\texttt{ConnectionLaplacian}); the table above can be regenerated by
invoking \verb|#print axioms| on each listed name. Continuous
fuzz-sweeps now span five rounds of the discovery swarm. The R3/R4
baseline exhausts $176{,}169$ connection-graph configurations
($n\le 5$ exhaustive; $n\in\{6,7,8\}$ wrap-subset / random-graph);
R5 extends this to a $456{,}950$-configuration sweep at $n\le 7$
targeting the cover-charpoly and kernel-dimension identities; R6
Stage-2 \texttt{FUZZER-A} (\texttt{findings/round6/stage2\_fuzzer\_A/})
certifies the M6 multiset-union identity at $1.5\cdot 10^{-15}$
relative precision on $3{,}727$ iso-configurations at $n=6$ plus a
random $n=7$ sample; R6 Stage-6 \texttt{FUZZER-B}
(\texttt{findings/round6/stage6\_fuzzer\_B/}) cross-checks the
shifted-determinant, trace-square, Frobenius-norm and
cross-polarisation identities of L17 at $n=7$ ($300$ graphs,
$14{,}011$ probes, zero counterexamples). Across all rounds, the
only non-integer numerical deviations are $20$ out of $24{,}971$
configurations at $n=8$ that show a $10^{-3}$-order relative error
in evaluating $\operatorname{charpoly}(\widetilde L^{\mathrm{cov}})$
through the companion matrix (attributable to condition-number
growth, bounded by standard float-polynomial perturbation theory;
max deviation $7.3\cdot 10^{-3}$). The underlying Lean identity
\texttt{cover\_charpoly\_eq\_scalar\_times\_mobius} is an exact
statement over $\ZZ[\lambda]$ and is unaffected. Fuzzer scripts and
JSON outputs live under
\texttt{findings/round\{3,4,5,6\}/}.

\paragraph{Counter-claims are also machine-shadowed.} Each
counter-example in \S\ref{sec:counter} (Examples~\ref{ex:naive-eq},
\ref{ex:bridge-wrap}, \ref{ex:cross-merge}, \ref{ex:contraction})
is emitted by the \texttt{NEGATOR} / \texttt{NEGATOR-FUZZY} passes
of the discovery swarm; the generating Python scripts live under
\path{findings/round3/negator/} and
\path{findings/round4/negator_fuzzy/} and can be re-run for any
$n\le 5$ (and $n=6$ with a wrap-subset cap) to regenerate the
counter-witnesses from scratch.

\begin{thebibliography}{9}

\bibitem{HararyCartwright}
F.~Harary.
\emph{On the notion of balance of a signed graph}.
Michigan Math.\ J.\ \textbf{2} (1953/54), 143--146.

\bibitem{Zaslavsky82}
T.~Zaslavsky.
\emph{Signed graphs}.
Discrete Appl.\ Math.\ \textbf{4} (1982), 47--74.

\bibitem{Chung}
F.~R.~K.~Chung.
\emph{Spectral Graph Theory},
CBMS Regional Conference Series in Mathematics~\textbf{92},
American Mathematical Society, Providence, 1997.

\bibitem{BiluLinial}
Y.~Bilu and N.~Linial.
\emph{Lifts, discrepancy and nearly optimal spectral gap}.
Combinatorica \textbf{26} (2006), 495--519.

\bibitem{Mathlib}
The Mathlib Community.
\emph{The Lean Mathematical Library}.
In: \emph{CPP~2020 Proceedings}, 367--381, 2020.

\bibitem{Lean4}
L.~de~Moura and S.~Ullrich.
\emph{The Lean~4 theorem prover and programming language}.
In: \emph{CADE~28 Proceedings}, LNCS~12699, 625--635, 2021.

\end{thebibliography}

\end{document}
